\clearpage
\VerdierPhysicalPageAnchor{0147}%
\begingroup
\small
\setlength{\parskip}{0pt}%
\setlength{\emergencystretch}{1.5em}%
\setlength{\abovedisplayskip}{3pt plus 1pt minus 1pt}%
\setlength{\belowdisplayskip}{3pt plus 1pt minus 1pt}%

\VerdierSemanticAnchor{ch2-2-3-4-d-complete}
\noindent これで d) が示された。次に f) を示す。
\VerdierSemanticAnchor{ch2-2-3-4-f}
$\mathcal D$ の対象 $Z$ が $Q^{-1}(\operatorname{Lib}_R(Q))$ に属するための
必要十分条件は、$Q$-右自由な対象 $Y$ と $\mathcal D/\mathcal B$ における
同型 $Q(Z)\xrightarrow[\sim]{m}Q(Y)$ が存在することである。
\VerdierSemanticAnchor{ch2-2-3-4-f-fraction}
この同型は $m=Q(s)^{-1}Q(t)$ と書ける。ここで
$s\in S_{\mathcal D}(\mathcal B)$、
$t\in\widehat{S_{\mathcal D}(\mathcal B)}$ である（(2.2.2) と
(2.2.4)）。よって $\mathcal D$ に次の図式がある：
\VerdierSemanticAnchor{ch2-2-3-4-f-diagram}
\begin{center}
 \begin{tikzcd}[column sep=4.1em,row sep=2.5em]
 Z \arrow[r,"t"] & X & Y \arrow[l,"s"']
 \end{tikzcd}\quad .
\end{center}
\VerdierSemanticAnchor{ch2-2-3-4-f-retraction}
a), iv) により、$s$ は $S_{\mathcal D}(\mathcal B)$ に属するレトラクション
$s'\colon X\longrightarrow Y$ をもつ (1.2.6)。したがって
$t'\colon Z\longrightarrow Y$、$t'\in\widehat{S_{\mathcal D}(\mathcal B)}$
が存在する。
\VerdierSemanticAnchor{ch2-2-3-4-f-converse}
逆にこのような射が存在すれば $Q(t')\colon Q(Z)\longrightarrow Q(Y)$ は
同型である。
\VerdierSemanticAnchor{ch2-2-3-4-g}
残る g) を示す。a), v) により双関手的同型
\VerdierSemanticAnchor{ch2-2-3-4-g-q1-adjunction}
\[
 \operatorname{Hom}_{Q^{-1}(\operatorname{Lib}_R(Q))}(Z,\Psi(W))
 \xrightarrow{\sim}
 \operatorname{Hom}_{\operatorname{Lib}_R(Q)}(Q_1(Z),W),
\]
があり、$\Psi$ は $Q_1$ の右随伴である。
\VerdierSemanticAnchor{ch2-2-3-4-g-inclusion-adjunction}
したがって包含関手を $i\colon L_R(Q)\lhook\joinrel\longrightarrow
Q^{-1}(\operatorname{Lib}_R(Q))$ と書けば、双関手的同型
\[
 \operatorname{Hom}_{Q^{-1}(\operatorname{Lib}_R(Q))}(Z,i(Y))
 \xrightarrow{\sim}
 \operatorname{Hom}_{L_R(Q)}(\Psi Q_1(Z),Y)\, .
\]
を得る。
\VerdierSemanticAnchor{ch2-2-3-4-g-proof-complete}
これは $\mathcal R=\Psi Q_1$ が包含関手の左随伴であることを示し、命題
(2.3.3) の証明を完了する。

\begingroup\itshape
\VerdierSemanticAnchor{ch2-2-3-5}
\medskip\noindent\textbf{命題 2.3.5.}\quad
\VerdierSemanticAnchor{ch2-2-3-5-setup}
$\mathcal D$ を三角圏、$\mathcal A$ と $\mathcal B$ を二つの充満三角部分圏
（$\mathcal B$ は厳密充満）とし、$S_{\mathcal D}(\mathcal B)$ と
$S_{\mathcal A}(\mathcal A\cap\mathcal B)$ をそれぞれ $\mathcal B$ と
$\mathcal A\cap\mathcal B$ に対応する $\mathcal D$ と $\mathcal A$ の
乗法的系とする。
\VerdierSemanticAnchor{ch2-2-3-5-a}
\par\noindent a) 次の条件は同値である：
\VerdierSemanticAnchor{ch2-2-3-5-a-i}
\par\noindent\hspace{1.2em}i) $X\in\operatorname{Ob}(\mathcal A)$、
$s\in S_{\mathcal D}(\mathcal B)$ である任意の射
$X\xrightarrow{s}X'$ は次の可換図式に組み込まれる：
\VerdierSemanticAnchor{ch2-2-3-5-a-i-diagram}
\begin{center}
 \begin{tikzcd}[column sep=4.1em,row sep=2.6em]
 X \arrow[r,"s"] \arrow[dr,"{s'}"'] & X' \arrow[d] &
 {s'\in S_{\mathcal A}(\mathcal A\cap\mathcal B)} \\
 & X'' &
 \end{tikzcd}\quad .
\end{center}
\endgroup
\endgroup
